% P1 of the public test corpus (docs/DESIGN.md#testing). Written by the project; dedicated to the public domain
% (CC0 1.0). A two-column article with the packages a physics paper typically loads. The figure file
% fig-response.pdf is generated by the test from code (test/engine/corpus.test.ts).
\documentclass[twocolumn,10pt]{article}
\usepackage[T1]{fontenc}
\usepackage[utf8]{inputenc}
\usepackage[letterpaper,margin=1.8cm]{geometry}
\usepackage{amsmath,amssymb}
\usepackage{graphicx}
\usepackage{booktabs}
\usepackage{siunitx}
\usepackage[numbers,sort&compress]{natbib}
\usepackage[hidelinks]{hyperref}

\title{A Damped Oscillator Measured Twice}
\author{The Test Corpus Project}
\date{}

\begin{document}
\maketitle

\begin{abstract}
We measure the decay of a damped harmonic oscillator with two independent methods and compare the fitted
quality factors. The two agree within \SI{2}{\percent}. The document exists to exercise a typical paper:
equations, a figure, a table, units and a bibliography.
\end{abstract}

\section{Introduction}
The damped harmonic oscillator is the standard model of a resonance~\citep{goldstein2002,landau1976}. Its
displacement obeys
\begin{equation}
  \ddot{x} + \frac{\omega_0}{Q}\,\dot{x} + \omega_0^2 x = 0,
  \label{eq:motion}
\end{equation}
where $\omega_0$ is the natural angular frequency and $Q$ the quality factor. For $Q \gg 1$ the amplitude
decays as $x(t) \propto e^{-\omega_0 t / 2Q}$, which is the basis of the ring-down method~\citep{knuth1984}.

\section{Methods}
We drove a mass of \SI{0.250}{\kilogram} on a spring with stiffness \SI{39.5}{\newton\per\metre}, which gives
$\omega_0 = \SI{12.57}{\radian\per\second}$. In the first method we fitted the ring-down envelope; in the
second we fitted the width $\Delta\omega$ of the driven response, using $Q = \omega_0 / \Delta\omega$.
Equation~\eqref{eq:motion} was integrated numerically to check both fits.

\begin{figure}[t]
  \centering
  \includegraphics[width=\linewidth]{fig-response}
  \caption{Driven response of the oscillator against drive frequency (arbitrary units).}
  \label{fig:response}
\end{figure}

\section{Results}
Figure~\ref{fig:response} shows the driven response, and Table~\ref{tab:q} lists the fitted quality factors.

\begin{table}[t]
  \centering
  \caption{Quality factors from the two methods.}
  \label{tab:q}
  \begin{tabular}{lS[table-format=3.1]S[table-format=1.1]}
    \toprule
    Method & {$Q$} & {Uncertainty} \\
    \midrule
    Ring-down       & 152.3 & 1.8 \\
    Resonance width & 149.6 & 2.4 \\
    \bottomrule
  \end{tabular}
\end{table}

\section{Discussion}
The two values differ by \SI{1.8}{\percent}, inside the combined uncertainty. Typesetting this comparison
needs nothing beyond the packages above~\citep{lamport1994,mittelbach2004}.

\bibliographystyle{unsrtnat}
\bibliography{paper}

\end{document}
